\documentclass[conference]{IEEEtran}

\usepackage{fontspec}
\usepackage{graphicx}
\usepackage{hyperref}
\usepackage{algpseudocode}
\usepackage{amsmath}
\usepackage{amssymb}

\algrenewcommand\algorithmiccomment[1]{$\triangleright$ #1}
\let\leq=\leqslant
\let\geq=\geqslant
\newcommand\MinP{\textsc{MinPlayer}}
\newcommand\MaxP{\textsc{MaxPlayer}}
\newcommand\DP{\textsc{DummiePlayer}}
\newcommand\SMPo{\textsc{SimpleMindedPlayer1}}
\newcommand\SMPt{\textsc{SimpleMindedPlayer2}}
\newcommand\GPA{\textsc{GeneralPurposeAdversary}}
\pagenumbering{arabic}

\title{%
\begin{center}
    \includegraphics[width=0.4\textwidth]{graphics/war.pdf}
\end{center}
\textbf{WAR: WAR with Auction Rounds} \\
{\large Report for \emph{Game Theory \& Networks} with \emph{Dr.~Dejun Yang}}
}
\author{
    Jack Rosenthal \\
    \texttt{jrosenth@mines.edu}
    \and Paul Sattizahn \\
    \texttt{psattiza@mines.edu}}
\date{2017-12-15}

\begin{document}

\maketitle
\thispagestyle{plain}
\pagestyle{plain}

\begin{abstract}
    War is a probabilistic card game commonly played in North America. In this
    paper, we investigate our own strategic variant of War called WAR (which
    stands for WAR with Auction Rounds) from a game theoretic approach, design
    strategies for playing the game, and analyze the effectiveness of the
    strategies using a computer simulation.
\end{abstract}

\section{Introduction}

\emph{War} is a card game played with two players and a single deck of cards.
The deck of cards is shuffled, and each player is dealt half of the deck. The
players repeatedly flip the top card of their deck, the player flipping the
higher valued card wins both cards and puts the cards in their \emph{wins
pile}. If the two players play the same card, they ``go to war''. In a war (the
resolution of a tie), both players discard the top three cards from their deck
to the pot, then play a fourth. The fourth card decides who wins the war. If
the fourth card is the same, the players go to war again until the war is
resolved. Once a player runs out of cards in their deck, they flip their wins
pile and continue to play. If they have no wins pile, the other player wins.

Studying War would be very boring from a game theoretical approach, as the game
involves no element of strategy. Instead, we propose a modified game: WAR -- a
recursive acronym standing for \textbf{WAR with Auction Rounds}. To distinguish
these two games, we will refer to the original game of War as \emph{Regular
War}, and our modified version as \emph{WAR} (as printed in all capital
letters).

WAR is a game played with $N$ players, each player starts with one suite of
cards (12 cards, labelled $2, 3, \ldots, 10, \mathrm{J}, \mathrm{Q},
\mathrm{K}$). The cards J, Q, and K have the values 11, 12, and 13,
respectively. The players then repeatedly \emph{choose} a card from
their hand to play, the cards are then flipped simultaneously and the player
who plays the highest card takes all the cards for their wins pile. Similar to
Regular War, when there is a tie for the highest card, all the players who
played the highest card go to war. Unlike Regular War, during a war, each of
the players involved get to \emph{choose} which three cards to discard and
which card to battle with. This process is repeated until the war is resolved.
Finally, just like Regular War, when a player is out of cards in their hand,
they pick up their wins pile if they have one, otherwise, they are out of the
game.

Normally, WAR is won by the last remaining player in the game. However, there
does exist the possibility that a game of WAR cannot be won by any of the
players. This will occur when all of the remaining players get in a tie, but
all players run out of cards before the tie is resolved.

To assist in your understanding of the game, the authors have produced a video
demonstration of the game. This video is available online at:
\begin{center}
\url{https://www.youtube.com/watch?v=KzWVCbNsjwc}
\end{center}

In this paper, we will define a model for the game, and discuss various
strategies (and their relation in terms of effectiveness to the other
strategies) we discovered.

\section{System Model}

WAR is a repeated static game; that is, all players move simultaneously, the
utility of one player depends on the decision of another player, and the state
of the game is not independent from round to round.

At the beginning of a game of WAR, every player begins with a known set of
cards, and after all of the events that occur in the game, the cards that are
exchanged are known to every other player. In theory, if every player were able
to count the cards that went by, they could determine what any other player has
in their hand and in their wins pile. We say that WAR is a game of \emph{perfect
information}, all players have complete knowledge of the state of the game
(with the exception of the information of the card another player is about to
play, which means our game differs slightly from the standard definition of
perfect information).

We will use the following mathematical notation in our strategies to describe
the system:
\begin{enumerate}
    \item $N$: The number of players in the game, where each player is
        identified with a number $n$, where $n \in \{1, \ldots{}, N\}$.
    \item $H_n$: The set of cards in player $n$'s hand.
    \item $W_n$: The set of cards in player $n$'s wins pile.
\end{enumerate}

In addition, we will define a utility function $u_n(H_n, W_n)$ which can not only
be used to evaluate the utility of a player at the beginning of a game, but
also evaluate how well a player is doing during the middle of a game.

\begin{displaymath}
    u_n(H_n, W_n) = \sum_{i=1}^{|H_n|}{H_{ni}} + \sum_{j=1}^{|W_n|}{W_{nj}}
\end{displaymath}

It follows that $u_n(H_n, W_n)$ for the winning player at the end of the game
is $N\sum_{i=2}^{13}{i} = 90N$ and the utility for all other players is $0$. In
the case when nobody wins, the utility for all players is $0$.

Figure~\ref{alg:war} formally defines a game of war using pseudocode.

\begin{figure}
    \begin{algorithmic}
        \Procedure{WAR}{$N$}
        \State \Comment{Initialize the hand and wins pile of all players}
        \For{$n$ in $\{1, \ldots, N\}$}
            \State $H_n \gets \{2, \ldots, 13\}$
            \State $W_n \gets \emptyset$
        \EndFor
        \State $z \gets \emptyset$
        \State \Comment{While there are still remaining players\dots}
        \While{$|\left\{H_n \,:\, n \in \left\{1, \ldots, N\right\} \,|\, |H_n| \geq 1\right\}| \geq 2$}
            \State $\mathbf{p} \gets \left[n \,:\, n \in \left\{1, \ldots, N\right\} \,|\, |H_n| \geq 1\right]$
            \State \Comment{Ask each player to select a card}
            \State $\mathbf{c} \gets []$
            \ForAll{$n$ in $\mathbf{p}$}
                \State $\mathbf{c}_n \gets$ \Call{SelectCard}{$n$, $H_n$}
            \EndFor
            \ForAll{$n$ in $\mathbf{p}$}
                \State Remove $\mathbf{c}_n$ from $H_n$ and reveal
            \EndFor
            \State $\mathbf{w} \gets [\mathbf{c}_n \,|\, \mathbf{c}_n \in \max(\mathbf{c})]$
            \State $z \gets z \cup \mathbf{c}$
            \If{$|\mathbf{w}| \ne 1$}
                \State \Comment{It's a tie. Go to war!}
                \While{$|\mathbf{w}| \leq 2$}
                    \State $\mathbf{c} \gets []$
                    \ForAll{$n$ in $\mathbf{w}$}
                        \If{$|H_n| \leq 3$}
                            \State $z \gets z \cup H_n$
                            \State $H_n \gets \emptyset$
                            \State Remove $n$ from $\mathbf{w}$ and continue
                        \EndIf
                        \State $d_1, d_2, d_3, \mathbf{c}_n \gets$ \Call{SelectTie}{$n$, $H_n$}
                        \State $z \gets z \cup \{d_1, d_2, d_3, \mathbf{c}_n\}$
                    \EndFor
                    \ForAll{$n$ in $\mathbf{c}$}
                        \State Remove $\mathbf{c}_n$ from $H_n$ and reveal
                    \EndFor
                    \State $\mathbf{w} \gets [\mathbf{c}_n \,|\, \mathbf{c}_n \in \max(\mathbf{c})]$
                \EndWhile
            \EndIf
            \If{$|\mathbf{w}| = 1$}
                \State $W_{\mathbf{w}_1} \gets W_{\mathbf{w}_1} \cup z$
                \State $z \gets \emptyset$
            \EndIf
            \ForAll{$n$ in $\mathbf{p}$}
                \If{$H_n = \emptyset$}
                    \State $H_n \gets W_n$
                    \State $W_n \gets \emptyset$
                \EndIf
            \EndFor
        \EndWhile
        \EndProcedure
    \end{algorithmic}
    \caption{A game of WAR, represented formally in pseudocode.}
    \label{alg:war}
\end{figure}

\section{Strategies}
To further analyze the game, we define strategies. Each strategy is assigned
to a type of player, such that all players of a given type will play WAR with
the corresponding strategy.

\subsection{\MinP}
The \MinP{} strategy will always pick its minimum card when
given a choice. In other words, for a \MinP{} $n$, $\textsc{SelectCard}(n,\,H_n) = \min(H_n)$.
In the case of a war (tie), the \MinP{} will discard
its lowest three cards, and then it selects its next lowest card to play.
In other words, for a \MinP{} $n$, $\textsc{SelectTie}(n,\,H_n) = (\min_1(H_n),\,\min_2(H_n),\,\min_3(H_n),\,\min_4(H_n))$.


\subsection{\MaxP}
The \MaxP{} strategy will always pick its maximum card when
given a choice. In other words, for a \MinP{} $n$, $\textsc{SelectCard}(n,\,H_n) = \max(H_n)$.
In the case of a war (tie), the \MaxP{} strategy will discard
similarly to a \MinP. The \MaxP{} discards its lowest three cards,
but it selects its highest card to play.
In other words, for a \MaxP{} $n$, $\textsc{SelectTie}(n,\,H_n) = (\min_1(H_n),\,\min_2(H_n),\,\min_3(H_n),\,\max(H_n))$.

\subsection{\DP}
The \DP{} implements a very simple strategy. Whenever the \DP{}
has the option between a set of choices, it will choose one of the choices
randomly, where each choice has an equal chance being picked. In other words,
for a \DP{} $n$, $\textsc{SelectCard}(n, H_n) = \textsc{RandomChoice}(H_n)$.
The cards to discard in the case of a tie are also picked randomly.
In other words, for a \DP{} $n$, $\textsc{SelectTie}(n,\,H_n) = \textsc{RandomSample}(4,\,H_n)$.


\subsection{\SMPo}
The \SMPo{} is the first strategy to take its opponents strategy into consideration when making a choice.
%TODO: Add stuff

\subsection{\SMPt}
The \SMPt{} is almost identical to the \SMPo, except it is {more,less} risky.
%TODO: Add stuff

\subsection{\GPA}
The \GPA{} is a strategy that takes advantage of other deterministic players'
strategies. When playing against a deterministic player, a \GPA{} player
will...
%TODO: Add stuff

\section{Simulation}

\section{Conclusion}


\end{document}
